% 本节由已运行的检验结果自动生成；原始结果位于第一问模型检验目录。
\section{第一问模型检验与鲁棒性分析}
\label{sec:q1-validation}

\subsection{检验对象与方法}
第一问属于确定性单日调度优化问题，不含需训练的预测器，故不使用随机五折交叉验证或分类一致性指标。随机划分时段还会破坏储能状态连续性。本文采用物理约束独立回代、连续松弛下界、无储能对照，以及扰动条件下的重优化和固定计划执行检验。

检验沿用原模型的单程充放电效率0.9、免费弃光、禁止售电和无额外电网接入容量限制假设。各情景初末储电量均为6000千瓦时，容量与充放电功率边界保持一致。

\subsection{物理有效性与最优性}
由原始负荷、光伏功率重新换算时段电量，独立计算母线平衡和储能递推残差：
\begin{align}
r_t^B&=g_t+R_t+d_t-L_t-c_t-w_t,\\
r_t^E&=E_t-E_{t-1}-\eta_{ch}c_t+d_t/\eta_{dis}.
\end{align}
两类最大绝对残差分别为 $2.27\times10^{-13}$ 和 $8.31\times10^{-13}$ 千瓦时，均远小于预设数值容差 $10^{-5}$ 千瓦时。母线平衡均方根残差为 $3.67\times10^{-14}$ 千瓦时，仅反映物理平衡精度，不代表预测误差。
充放电功率、储能容量、非负性和弃光上界的最大违反量为零，日末储能偏差为零，同时充放电时段数为零。另检查
\begin{equation}
\sum_tg_t=\sum_t(L_t-R_t+w_t)+(1-\eta_{ch})\sum_tc_t
+(1/\eta_{dis}-1)\sum_td_t,
\end{equation}
其残差为 $5.46\times10^{-12}$ 千瓦时。原保存调度也通过回代，新旧购电成本一致。

独立建立去除充放电互斥的连续模型，保留其余物理约束，得到原问题最优成本下界。无储能成本为 48052.04659085 元，连续松弛下界为 35126.94858929 元，最终两阶段成本为 35126.94859929 元，节费率为 26.8981\%。最终成本与下界仅差 $10^{-5}$ 元，对应第二阶段减少吞吐量时允许的购电费容差，小于 $10^{-4}$ 元的检验阈值，支持数值容差内的全局最优性。物理可行性与下界证据共同表明，模型对给定数据和假设有效，但尚非实际预测准确性或所有工况适用性的证明。

\subsection{输入扰动与重优化}
设扰动后的负荷及光伏为
\begin{equation}
\widetilde P_t^L=P_t^L(1+\alpha\xi_t^L),\qquad
\widetilde P_t^{PV}=P_t^{PV}(1+\alpha\xi_t^{PV}),\qquad
\xi\sim U[-1,1].
\end{equation}
取 $\alpha=5\%,10\%,20\%$，其中10\%表示最大相对扰动幅度，非标准差；电价保持不变，夜间零光伏保持为零。逐时段独立组每时段抽样，一小时相关块组每六个连续时段共享误差，负荷与光伏分别独立抽样。随机种子为20260911，每种结构和幅度生成100个情景，共600个；同一结构不同幅度共享基础随机数。

成本变化率定义为 $100(C_\omega^*-C^*)/C^*$。表中区间为100个情景的2.5\%与97.5\%经验分位区间，并非均值置信区间或真实运行保证。
\begin{table}[htbp]
\centering\small
\caption{随机扰动下重优化成本的相对变化}
\begin{tabular}{lrrrr}
\toprule
结构 & 幅度 & 均值 & 经验分位区间 & 最大绝对变化\\
\midrule
逐时段独立 & 5\% & +0.031\% & [-0.979\%, +1.136\%] & 1.506\%\\
逐时段独立 & 10\% & +0.111\% & [-1.936\%, +2.339\%] & 3.174\%\\
逐时段独立 & 20\% & +0.416\% & [-3.723\%, +5.055\%] & 6.775\%\\
一小时相关块 & 5\% & +0.131\% & [-2.997\%, +3.485\%] & 5.854\%\\
一小时相关块 & 10\% & +0.347\% & [-5.910\%, +7.097\%] & 11.412\%\\
一小时相关块 & 20\% & +1.096\% & [-11.535\%, +14.763\%] & 21.382\%\\
\bottomrule
\end{tabular}
\end{table}
600个情景均通过求解和独立回代，最大违反量为 $1.66\times10^{-11}$ 千瓦时。
10\%独立扰动和一小时相关块扰动下，最大绝对成本变化分别为 3.174\% 与 11.412\%。各情景均根据完整扰动路径重新求解，因而属于完全信息下的敏感性试验，不能替代日前固定计划的执行检验。

\subsection{原计划执行能力}
首先将原购电、充放电和弃光全部冻结，按扰动数据回代平衡，将负残差累加为待补救缺电量，正残差记录为富余量；此处缺电量并非已经模拟真实停电。
其次只固定电池充放电，允许按原电价即时调整购电与弃光，令
\begin{align}
\widetilde g_t&=[\widetilde L_t-\widetilde R_t+c_t^0-d_t^0]_+,\\
\widetilde w_t&=[-(\widetilde L_t-\widetilde R_t+c_t^0-d_t^0)]_+.
\end{align}
还需检查 $\widetilde w_t\le\widetilde R_t$；否则在禁止售电条件下原电池放电不可吸收。仅对可行情景计算后悔率 $100(C_\omega^{fix}-C_\omega^*)/C_\omega^*$，不可行情景成本和后悔率留空。
\begin{table}[htbp]
\centering\small
\caption{固定计划的执行检验}
\begin{tabular}{lrrrr}
\toprule
结构 & 幅度 & 平均缺电量 & 电池固定可行率 & 平均后悔率\\
\midrule
逐时段独立 & 5\% & 1629.311 & 0\% & 不适用\\
逐时段独立 & 10\% & 3258.623 & 0\% & 不适用\\
逐时段独立 & 20\% & 6517.246 & 0\% & 不适用\\
一小时相关块 & 5\% & 1662.992 & 13\% & 1.559\%\\
一小时相关块 & 10\% & 3325.983 & 13\% & 3.011\%\\
一小时相关块 & 20\% & 6651.967 & 9\% & 5.570\%\\
\bottomrule
\end{tabular}
\end{table}
表中缺电量单位为千瓦时；可行率对应固定电池、允许电网补救的条件，后悔率只统计可行情景。本补救假设未计应急溢价或调整惩罚，不能作为后续小问的结算结果。电池调度相对变化另按充放电向量绝对差总和除以基准吞吐量记录于明细。同价时段存在等价最优解时，调度向量变化不宜直接解释为经济性不稳定。

10\%独立扰动下，固定电池方案在100个情景中均不可行；一小时相关块扰动下仅13个可行。某些时段负荷下降导致原定电池放电超过可吸收负荷，在禁止售电、禁止任意丢弃电池放电时，单纯调整购电无法恢复平衡。必须允许电池调度随实际输入调整。该发现针对本次返回的特定最优调度，并非对所有等价最优解的共同证明。

\subsection{系统性偏差与效率假设}
\begin{table}[htbp]
\centering\small
\caption{系统性输入偏差的压力检验}
\begin{tabular}{lrrr}
\toprule
情景 & 成本（元） & 成本变化 & 节费率\\
\midrule
负荷增加百分之十 & 42745.8086 & +21.690\% & 22.276\%\\
光伏减少百分之十 & 38422.5611 & +9.382\% & 23.496\%\\
负荷增加且光伏减少百分之十 & 46572.7646 & +32.584\% & 19.488\%\\
负荷减少且光伏增加百分之十 & 26509.7370 & -24.532\% & 34.107\%\\
电价整体增加百分之十 & 38639.6435 & +10.000\% & 26.898\%\\
电价整体减少百分之十 & 31614.2537 & -10.000\% & 26.898\%\\
\bottomrule
\end{tabular}
\end{table}
负荷全天增加10\%且光伏全天减少10\%时，成本变化为 +32.584\%，说明持续偏差无法像零均值独立误差一样相互抵消。电价整体同比变化时，购电成本同幅度变化，也支持目标函数量纲正确。
\begin{table}[htbp]
\centering\small
\caption{效率解释及参数敏感性}
\begin{tabular}{lrrr}
\toprule
单程效率 & 往返效率 & 成本（元） & 相对主方案变化\\
\midrule
0.800000 & 0.6400 & 38223.6516 & +8.816\%\\
0.850000 & 0.7225 & 36603.4065 & +4.203\%\\
0.900000 & 0.8100 & 35126.9486 & +0.000\%\\
0.948683 & 0.9000 & 33801.4956 & -3.773\%\\
0.950000 & 0.9025 & 33767.0326 & -3.871\%\\
1.000000 & 1.0000 & 32377.0883 & -7.828\%\\
\bottomrule
\end{tabular}
\end{table}
若题设90\%为往返效率，则应取 $\eta_{ch}=\eta_{dis}=\sqrt{0.9}$，成本为 33801.4956 元，较主方案变化 -3.773\%。主口径应根据题意或设备定义确定，不能以低费用作为选择依据。

\subsection{改进方向与结论边界}
\begin{enumerate}
\item 对持续误差和系统性偏差，使用真实预测残差校准偏差、尾部及相关性，引入连续时间块或联合情景、滚动优化及储能安全余量，再在相同信息集下与本模型比较。
\item 原计划遇到供需缺口时，应明确调整时点与费用，加入实际结算下的应急购电与调整变量。若存在电网接入上限，增加 $g_t\le\overline P^g\Delta t$，并独立统计缺供量和缺供时段。本次未验证有限接入容量下的适应性。
\item 设备数据若不支持常效率假设，可采用随功率、储电量或温度变化的分段线性效率，并标定循环损耗费用。当前第二层减少吞吐量仅筛选近似同成本解，不能替代完整衰减模型。
\item 若允许售电或存在弃光惩罚，应增加收益、惩罚和容量限制，必要时增加购售电互斥。本次单日无弃光不代表该假设对其他情景均无影响。
\item 扩展至多季节、多典型日真实样本外预测驱动的调度检验，报告成本、缺供、应急购电及尾部风险。若单独建立预测模块，使用按时间推进的验证；不能将本次合成扰动等同于全年泛化检验。
\end{enumerate}
本次物理残差已达浮点精度量级，优化有效性得到数值支持。抗扰结论限于本次扰动分布及补救假设；重优化成本稳定不等于原计划可直接执行。优先改进应是持续误差建模、执行补救机制和效率、结算口径的核实。
独立扰动下重优化成本较稳定，但持续性误差与系统性偏差显著增加成本波动，原固定电池计划的执行能力不足，故不将现模型概括为具有很强抗干扰能力。

完整检验代码为项目根目录的\texttt{problem1\_validate.py}，运行后由\texttt{working/write\_problem1\_validation\_report.py}生成本节；逐情景结果、随机种子、版本与输入校验值均保存在第一问模型检验结果目录。本文尚未完成排版编译与页面验收。
